% Generated from locale/tr/content/set-theory/spine/stagesbasics.tex; do not hand-edit.


\olfileid[tr]{sth}{spine}{Valphabasic}
\olsection{Aşamaların Temel Özellikleri}

$V_\alpha$'ların tanımının temellendirmedeki önemini ortaya çıkarmak için
onlara ilişkin birkaç temel sonuç sunacağız. Bir tanımla başlayalım:
\footnote{``Kuvvetli'' için standart bir terminoloji yoktur; bu,
\citet{ButtonLT1} tarafından kullanılan addır.}
\begin{defn}
	$A$ kümesi, $\forall x((\exists y \in A)x \subseteq y \lif x \in A)$
	ise \emph{kuvvetlidir}.
\end{defn}
\begin{lem}\ollabel{Valphabasicprops}
Her $\alpha$ sıra sayısı için:
\begin{enumerate}
	\item\ollabel{Valphatrans} Her $V_\alpha$ geçişlidir.
	\item\ollabel{Valphapotent} Her $V_\alpha$ kuvvetlidir.
	\item\ollabel{Valphacum} $\gamma \in \alpha$ ise $V_\gamma
	\in V_\alpha$ olur (ve dolayısıyla \olref{Valphatrans} gereği ayrıca
	$V_\gamma\subseteq V_\alpha$ olur).
\end{enumerate}
\end{lem}

\begin{proof}
Bunu (eşzamanlı) sonlu ötesi tümevarımla ispatlayacağız. Tümevarım için
\olref{Valphatrans}--\olref{Valphacum} sonuçlarının her $\beta<\alpha$
sıra sayısı için geçerli olduğunu varsayın.

$\alpha=\emptyset$ durumu açıktır.

$\alpha=\ordsucc{\beta}$ olsun. \olref{Valphacum} sonucunu göstermek için
$\gamma\in\alpha$ ise varsayım gereği $V_\gamma\subseteq V_\beta$ ve
dolayısıyla $V_\gamma\in\Pow{V_\beta}=V_\alpha$ olur.
\olref{Valphapotent} sonucunu göstermek için
$A\subseteq B\in V_\alpha$, yani $A\subseteq B\subseteq V_\beta$
olduğunu varsayın; o zaman $A\subseteq V_\beta$ ve dolayısıyla
$A\in V_\alpha$ olur. \olref{Valphatrans} sonucunu göstermek için
$x\in A\in V_\alpha$ ise $A\subseteq V_\beta$ ve dolayısıyla
$x\in V_\beta$ olduğunu gözlemleyin. Varsayım gereği $V_\beta$ geçişli
olduğundan $x\subseteq V_\beta$ ve böylece $x\in V_\alpha$ olur.

$\alpha$ bir limit sıra sayısı olsun. \olref{Valphacum} sonucunu göstermek
için $\gamma\in\alpha$ ise
$\gamma\in\ordsucc{\gamma}\in\alpha$ olur; dolayısıyla varsayım gereği
$V_\gamma\in V_{\ordsucc{\gamma}}$ ve böylece
$V_\gamma\in\bigcup_{\beta\in\alpha}V_\beta=V_\alpha$ olur.
\olref{Valphatrans} ile \olref{Valphapotent} için yalnızca geçişli
(sırasıyla kuvvetli) kümelerin birleşiminin geçişli (sırasıyla kuvvetli)
olduğunu gözlemleyin.
\end{proof}

\begin{lem}\ollabel{Valphanotref}
Her $\alpha$ sıra sayısı için $V_\alpha\notin V_\alpha$.
\end{lem}

\begin{proof}
Sonlu ötesi tümevarımla. Açıktır ki
$V_\emptyset\notin V_\emptyset$.

$V_{\ordsucc{\alpha}}\in V_{\ordsucc{\alpha}}=\Pow{V_\alpha}$ ise
$V_{\ordsucc{\alpha}}\subseteq V_\alpha$ olur; ayrıca
\olref{Valphabasicprops} gereği
$V_\alpha\in V_{\ordsucc{\alpha}}$ olduğundan $V_\alpha\in V_\alpha$
elde edilir. Karşıt-ters olarak $V_\alpha\notin V_\alpha$ ise
$V_{\ordsucc{\alpha}}\notin V_{\ordsucc{\alpha}}$ olur.

$\alpha$ bir limit ve
$V_\alpha\in V_\alpha=\bigcup_{\beta\in\alpha}V_\beta$ ise bir
$\beta\in\alpha$ için $V_\alpha\in V_\beta$ olur. Fakat o zaman ayrıca
$V_\beta\in V_\alpha$ olur; böylece
\olref{Valphabasicprops} iki kez kullanıldığında
$V_\beta\in V_\beta$ elde edilir. Karşıt-ters olarak bütün
$\beta\in\alpha$ için $V_\beta\notin V_\beta$ ise
$V_\alpha\notin V_\alpha$ olur.
\end{proof}

\begin{cor}
Her $\alpha,\beta$ sıra sayısı için $\alpha\in\beta$ ancak ve ancak
$V_\alpha\in V_\beta$.
\end{cor}

\begin{proof}
\Olref{Valphabasicprops} bir yönü verir. Tersine
$V_\alpha\in V_\beta$ olduğunu varsayın. \olref{Valphanotref} gereği
$\alpha\neq\beta$ olur. Ayrıca $\beta\notin\alpha$'dır; yoksa
$V_\beta\in V_\alpha$ ve dolayısıyla
\olref{Valphabasicprops} iki kez kullanıldığında
$V_\beta\in V_\beta$ olurdu; bu da \olref{Valphanotref} ile çelişir.
Öyleyse Üçlü Karşılaştırma gereği $\alpha\in\beta$ olur.
\end{proof}
\noindent
Bütün bunlar her $V_\alpha$'yı hiyerarşinin $\alpha$'ıncı aşaması olarak
düşünmemizi sağlar. Nedenini görelim.

Sıra sayıları kullanmamız yineleme kavramını izlediğinden, $V_\alpha$'ların
\emph{yinelemeli} bir süreçte oluşturulduğunu düşünebiliriz. Ayrıca bir
aşama ötekinden önce oluşturulmuşsa, yani $V_\beta\in V_\alpha$, yani
$\beta\in\alpha$ ise, $V_\beta\subseteq V_\alpha$ olduğundan oluşturma
sürecimiz \emph{birikimlidir}. Son olarak gerçekten de daha önceki herhangi
bir aşamada kullanılabilir olan kümelerin \emph{bütün} olası topluluklarını
oluştururuz; çünkü her ardıl aşama $V_{\ordsucc{\alpha}}$, önceli
$V_\alpha$'nın kuvvet kümesidir.

Kısacası $\ZFminus$ ile birikimli-yinelemeli küme hiyerarşisi görüşümüzü
dile getirme işini \emph{neredeyse} tamamladık. (Elbette Yerine Koyma'yı
hâlâ gerekçelendirmemiz gerekir.)

